%! AP Statistics — Unit 3 Cover Sheet
\documentclass[10pt]{article}
\usepackage{apstats-article}
\usepackage{apstats-boxes}

%=============================================================================
\begin{document}

% ── Full-bleed navy banner ────────────────────────────────────────────────────
\begin{tikzpicture}[remember picture, overlay]
  \fill[navy] (current page.north west)
    rectangle ([yshift=-1.4in]current page.north east);
\end{tikzpicture}
\vspace{-1.3in}

\begin{center}
  {\color{white}\bfseries\fontsize{28}{34}\selectfont AP Statistics}\\[8pt]
  {\color{goldacc}\bfseries\Large Shepherd \quad 2026--2027}\\[14pt]
  {\color{white}\bfseries\LARGE Unit 3: Collecting Data}\\[4pt]
\end{center}

\vspace{0.30in}

% ── Unit overview ─────────────────────────────────────────────────────────────
\begin{tcolorbox}[colback=sky,colframe=navy,boxrule=0.9pt,arc=2mm,
    left=4mm,right=4mm,top=3mm,bottom=3mm]
  \textbf{Unit Overview.}\quad
  Unit 3 examines how data are collected and why the collection method determines
  what conclusions can validly be drawn.  Students begin by questioning whether data
  collected without chance can be trusted, then distinguish observational studies from
  experiments.  The sampling half of the unit covers simple random, stratified, cluster,
  and systematic samples, and identifies the bias sources that arise when chance is
  absent.  The experimental half introduces the components and principles of
  well-designed experiments --- random assignment, control, replication, blinding, and
  blocking --- and closes with the logic of statistical inference: how random assignment
  allows researchers to attribute statistically significant differences to the treatments.
\end{tcolorbox}

\vspace{0.22in}

% ── Lessons in this unit ──────────────────────────────────────────────────────
\renewcommand{\arraystretch}{1.6}
\begin{skillbox}[Lessons in This Unit]{goldbox}
\begin{tabularx}{\linewidth}{@{} c >{\bfseries}p{0.44\linewidth} X @{}}
  \textbf{\#} & \textbf{Title} & \textbf{Focus} \\
  \hline
  3.1 & Introducing Statistics: Do the Data We Collected Tell the Truth?
    & Questioning collection methods; chance vs.\ non-chance data. \\
  \hline
  3.2 & Introduction to Planning a Study
    & Population vs.\ sample; observational study vs.\ experiment; generalization. \\
  \hline
  3.3 & Random Sampling and Data Collection
    & SRS, stratified, cluster, systematic, and census; choosing a method. \\
  \hline
  3.4 & Potential Problems with Sampling
    & Voluntary response, undercoverage, nonresponse, and question-wording bias. \\
  \hline
  3.5 & Introduction to Experimental Design
    & Components of experiments; principles of well-designed experiments; blocking. \\
  \hline
  3.6 & Selecting an Experimental Design
    & Choosing among completely randomized, block, and matched-pairs designs. \\
  \hline
  3.7 & Inference and Experiments
    & Statistically significant results; causal conclusions; scope of inference. \\
  \hline
\end{tabularx}
\end{skillbox}

\vspace{0.18in}

% ── Standards covered ─────────────────────────────────────────────────────────
\begin{skillbox}[AP Statistics Big Ideas \& Learning Objectives --- Unit 3]{greenbox}
\begin{tabularx}{\linewidth}{@{} >{\bfseries}p{0.14\linewidth} X @{}}
  VAR-1.E & Identify questions to be answered about data collection methods;
            recognize that non-chance methods result in untrustworthy conclusions. \\
  DAT-2.A & Identify the type of study (observational vs.\ experiment); distinguish
            population from sample. \\
  DAT-2.B & Identify appropriate generalizations from observational studies;
            recognize that causation cannot be established from observational data. \\
  DAT-2.C & Identify and describe sampling methods: SRS, stratified, cluster,
            systematic, and census. \\
  DAT-2.D & Explain why a particular sampling method is or is not appropriate for
            a given situation. \\
  DAT-2.E & Identify potential sources of bias: voluntary response, undercoverage,
            nonresponse, and response bias. \\
  VAR-3.A & Identify the components of an experiment: experimental units, factors,
            treatments, response variable, and confounding variables. \\
  VAR-3.B & Describe the elements of a well-designed experiment: comparison, random
            assignment, replication, and control of confounding variables. \\
  VAR-3.C & Compare experimental designs: completely randomized, single- and
            double-blind, control group, block design, and matched pairs. \\
  VAR-3.D & Explain why a particular experimental design is appropriate. \\
  VAR-3.E & Interpret results of a well-designed experiment; determine whether
            differences are statistically significant and what scope of inference applies. \\
\end{tabularx}
\end{skillbox}

\end{document}
